\section{Discussion}
\label{sec:discussion}

\para{The E1/E2 asymmetry.} The two sycophancy probes ask the model to fail in different directions. \eone (``are you sure?'') asks the model to \emph{defend a correct answer} when challenged -- the right behavior is a refusal to change. \etwo (productive friction) asks the model to \emph{volunteer a correction} the user did not ask for -- the right behavior is unsolicited disagreement. The first kind of failure has been the explicit target of safety training for two years; the second has not. We read the asymmetric pattern in our data (small persona effects on \eone, large persona effects on \etwo) as evidence that frontier-model RLHF has solved the easy direction of sycophancy (resisting ``are you sure?'') but not the hard direction (volunteering correction). The hard direction is the one that matters for collaboration -- and it is the one where persona moves the needle.

\para{Refuting a strong form of the Personality Illusion.} \citet{han2025personality} showed that persona injection moves Big-Five self-reports (\(\beta=3.95\)) but does not move behavior on a sycophancy task (\(\beta=0.03\), \(p=0.67\)). We replicate their null for demographic personas on \gptmini (\pdemo \etwo sycophancy \(\Delta = +1.3\)~\pp, \(p=0.56\)) -- the persona class their work tested. But we find that a \emph{belief-based} persona moves \etwo sycophancy by \(18.75\)--\(21.25\) percentage points (\(p < 0.01\) on both models), alongside a Cohen's \(d > 1.3\) behavioral signature at the embedding level (\ethree). The natural reading is that Han et al.'s finding is real for the persona families their work tested (Big-Five and role-based), but does not generalize to personas built around explicit epistemic commitments. The distinct-persona hypothesis predicts exactly this gap; the gap is large.

\para{What kinds of conversations need a believing other?} The submitter's framing posed this as the open question. Our data point to a partial taxonomy. In conversations where the user can verify on their own (and probably will -- \eone is essentially a quiz the user can check), modern safety training has already done much of the work and persona contributes little. In conversations where the user has stated a confident falsity and is asking the model to elaborate (\etwo -- they think they already know), the user will \emph{not} verify, and the only thing standing between them and a misinformation cascade is the model's willingness to push back. This is exactly the regime in which our largest effects appear. The provisional answer to the submitter's question is: \emph{conversations where the user has already committed to a wrong belief and will not check require an interlocutor with stance.}

\para{``Helpful assistant'' is not neutral.} The biggest single anomaly in our data is that the most common system prompt in industry deployment -- ``You are a helpful assistant.'' -- nearly quadruples the \eone fold rate on \gptone, from \(7.1\%\) to \(27.3\%\) (\(p=0.014\)), relative to no system prompt. The literal default is therefore not a neutral prior; it is an \emph{actively sycophantic} prior on this model. This deserves replication and extension before it can be taken as a generalization, but if it holds it has direct deployment consequences: the most common framing in industry is making the failure mode it is trying to mitigate worse, not better.

\para{Mechanism speculation.} Why would a system prompt that simply names an epistemic stance produce a \(\sim 20\)~\pp behavior change when one that names a person (\pdemo) does not? One possible mechanism is that the belief persona \emph{rephrases the helpfulness objective}: ``be helpful'' under the default RLHF preference distribution often resolves to ``be agreeable,'' while ``be helpful given that pressure is not evidence'' relocates the local optimum to a different point in response-space. The \ethree result -- that response style shifts as much as behavior does -- is consistent with this: the model is generating from a different conditional distribution, not merely truncating the same one.

\para{Limitations.} We are explicit about six.
\begin{itemize}[leftmargin=*,itemsep=0pt,topsep=0pt]
    \item \textbf{One vendor.} Both models are OpenAI \texttt{gpt-4.1}-family. The effect could be specific to OpenAI's system-prompt fine-tuning. Cross-vendor replication on Claude Sonnet-4.5, Gemini-2.5, and Llama-3.x is necessary before claiming generality.
    \item \textbf{LLM-as-judge for \etwo.} The \etwo correction/sycophancy labels come from \gptmini, blind to condition and grounded by a TruthfulQA correct answer. We did not run a human-validation pass. A 20\% spot check would strengthen the claim; the literature flags judge bias as a real risk.
    \item \textbf{Single persona instantiation.} We test \emph{one} belief-persona text. The effect could depend on specific wording rather than the category. A persona-text sensitivity analysis with three variants each of \pcai and \pbelief would separate category effect from wording effect.
    \item \textbf{\(n=80\) per cell.} Sufficient for the \(\sim 20\)~\pp effects we find but underpowered for smaller effects (\eg \pbelief vs \pnone on \eone \gptmini is in the predicted direction but not significant at this \(n\)).
    \item \textbf{TruthfulQA training contamination.} \truthfulqa is partly in training data, which biases against finding any persona effect (the model would know the truth regardless). The fact that a \(\sim 20\)~\pp effect appears \emph{despite} this is conservative.
    \item \textbf{Single-turn only.} \citet{liu2025truthdecay} show sycophancy compounds over multi-turn dialog. Whether the belief persona holds up at 3 or 7 turns is an open empirical question.
\end{itemize}

\para{Broader implications.} If the effect generalizes, the practical implication is that any deployment that needs principled pushback -- coding pair, proofreader, devil's advocate, fact-checker, tutor -- should consider replacing the default ``You are a helpful assistant'' framing with an explicit epistemic stance. The negative finding from \citet{han2025personality} should be read narrowly: persona is not action-controlling \emph{when persona is demographic or trait-based}. It is action-controlling when the persona names commitments about evidence and pressure. The framing of identity matters.
